\section{相关工作}

\subsection{医学检索增强问答}

近期基准研究已证实检索质量主导下游医学问答性能——其影响往往超过生成模型的选择~\citep{xiong2024benchmarking}。MedRAG/MIRAGE 基准系统评估了语料、检索器和主干模型，发现即便强生成模型在接收弱证据时也会显著退化。迭代式医学 RAG 通过生成后续检索查询来应对这一问题~\citep{xiong2024imedrag}，但其核心证据选择步骤仍是成对的——每个段落独立于查询进行评分。这些系统共同关注的是\emph{使用}检索证据进行生成；它们共同缺乏的是利用连接段落之间以及段落与结构化知识之间的高阶生物医学关系来重排序证据的机制。

\subsection{生物医学问答数据集与语义检索}

BioASQ~\citep{tsatsaronis2015bioasq} 及其人工整理的问答扩展 BioASQ-QA~\citep{krithara2023bioasqqa} 是生物医学证据检索的主要评估框架，将问题与金标准段落级标注配对。PubMedQA~\citep{jin2019pubmedqa} 和 BEIR 的 NFCorpus~\citep{thakur2021beir} 提供了不同证据结构的补充测试平台。在生物医学语义检索方面，MedCPT~\citep{jin2023medcpt} 在大规模 PubMed 搜索日志上训练双编码器检索器和交叉编码器重排序器，建立了强神经基线。然而，这些语义模型孤立地操作查询-段落对。它们不利用任何结构化生物医学知识——MeSH 层级、生物医学实体图或精选关系数据库——尽管这些资源在生物医学领域是免费可得的。这一缺口驱动了我们的工作：将结构化生物医学知识加入检索特征能否进一步提升证据重排序？

\subsection{图与超图 RAG 方法}

越来越多的工作将图结构应用于检索信息。GraphRAG~\citep{edge2024graphrag} 将段落组织为实体-关系图并执行局部到全局搜索用于摘要。Medical Graph RAG~\citep{wu2024medgraphrag} 通过将用户文档与可信来源和受控词表连接来建立医学回答的证据基础。超图扩展主张 n 元关系——同时连接两个以上节点——更适合表示复杂信息结构。HyperGraphRAG~\citep{luo2025hypergraphrag} 使用超边进行段落聚类，跨粒度 HGRAG~\citep{wang2026hgrag} 在不同检索粒度上建模实体和段落。Medical HyperRAG~\citep{xu2026medicalhyperrag} 构建 ClinBridge 超图组织病历、病例和指南用于患者中心问答。

这些方法有一条共同的线索：它们在\emph{缺乏领域概念标注}的情况下应用图结构。边通常由词汇重叠、嵌入相似度或共现统计构建。我们的工作检验了以下假设：这是不够的——原始拓扑在没有医学知识约束的情况下可能引入排序噪声而非信号。通过在同一超图架构上有无 MeSH、实体和 PrimeKG 标注的系统比较，我们提供了直接证据表明领域知识约束是结构特征在生物医学证据重排序中有用的必要条件。

\subsection{受控生物医学知识}

MeSH~\citep{nlm2026mesh} 是由美国国家医学图书馆维护的层级化受控词表，为生物医学文献索引提供标准化标注。PrimeKG~\citep{chandak2023primekg} 整合了跨疾病、蛋白、药物和表型的多模态精准医学关系。两种资源均可公开获取，是生物医学信息检索系统的基础，但此前均未被系统性地集成到证据检索的学习排序流程中。我们的框架以 MeSH 层级特征作为主要结构化知识信号——占证据救回的 75.9\%——并将 PrimeKG 作为辅助来源，其精确名称覆盖限制了独立贡献。
